\documentclass[11pt,a4paper]{article}

% ── Packages ──────────────────────────────────────────────────────────────────
\usepackage[a4paper, top=22mm, bottom=22mm, left=22mm, right=22mm]{geometry}
\usepackage[T1]{fontenc}
\usepackage[utf8]{inputenc}
\setlength{\headheight}{22pt}
\usepackage{xcolor}
\usepackage{tcolorbox}
\usepackage{listings}
\usepackage{booktabs}
\usepackage{array}
\usepackage{tabularx}
\usepackage{enumitem}
\usepackage{titlesec}
\usepackage{fancyhdr}
\usepackage{parskip}
\usepackage{mdframed}
\usepackage{amsmath}
\usepackage{amssymb}
\usepackage{graphicx}
\usepackage{hyperref}
\usepackage{tikz}
\usetikzlibrary{shapes.geometric, arrows.meta, positioning, fit, backgrounds}

\tcbuselibrary{skins, breakable, listings}

% ── Colour palette ────────────────────────────────────────────────────────────
\definecolor{primaryBlue}{HTML}{1D4ED8}
\definecolor{lightBlue}{HTML}{DBEAFE}
\definecolor{accentPurple}{HTML}{6D28D9}
\definecolor{lightPurple}{HTML}{EDE9FE}
\definecolor{codeGreen}{HTML}{15803D}
\definecolor{codeBg}{HTML}{F1F5F9}
\definecolor{alertRed}{HTML}{B91C1C}
\definecolor{lightRed}{HTML}{FEE2E2}
\definecolor{amber}{HTML}{B45309}
\definecolor{lightAmber}{HTML}{FEF3C7}
\definecolor{teal}{HTML}{0F766E}
\definecolor{lightTeal}{HTML}{CCFBF1}
\definecolor{darkText}{HTML}{1E293B}
\definecolor{mutedText}{HTML}{64748B}
\definecolor{borderGray}{HTML}{CBD5E1}
\definecolor{sectionBar}{HTML}{3B82F6}
\definecolor{headerBg}{HTML}{1E3A5F}
\definecolor{noteYellow}{HTML}{FFFBEB}
\definecolor{noteYellowBorder}{HTML}{F59E0B}

% ── Page style ────────────────────────────────────────────────────────────────
\pagestyle{fancy}
\fancyhf{}
\renewcommand{\headrulewidth}{0.4pt}
\renewcommand{\headrule}{\hbox to\headwidth{\color{sectionBar}\leaders\hrule height \headrulewidth\hfill}}
\fancyhead[L]{\small\color{mutedText}\textbf{Chat Templates in Transformers}}
\fancyhead[R]{\small\color{mutedText}apply\_chat\_template() — Complete Notes}
\fancyfoot[C]{\small\color{mutedText}\thepage}

% ── Section formatting ────────────────────────────────────────────────────────
\titleformat{\section}
  {\color{primaryBlue}\large\bfseries}
  {}
  {0em}
  {\colorbox{lightBlue}{\parbox{\dimexpr\linewidth-2\fboxsep\relax}{\color{primaryBlue}\thesection\quad #1}}}
  [\vspace{2pt}]

\titleformat{\subsection}
  {\color{accentPurple}\normalsize\bfseries}
  {\thesubsection}
  {0.5em}
  {}
  [\vspace{1pt}\color{accentPurple}\hrule height 0.5pt\vspace{3pt}]

\titlespacing*{\section}{0pt}{14pt}{6pt}
\titlespacing*{\subsection}{0pt}{10pt}{4pt}

% ── Code listing style ────────────────────────────────────────────────────────
\lstdefinestyle{pycode}{
  language=Python,
  backgroundcolor=\color{codeBg},
  basicstyle=\ttfamily\small,
  keywordstyle=\color{primaryBlue}\bfseries,
  stringstyle=\color{codeGreen},
  commentstyle=\color{mutedText}\itshape,
  numberstyle=\tiny\color{mutedText},
  numbers=left,
  numbersep=8pt,
  frame=l,
  framesep=4pt,
  rulecolor=\color{sectionBar},
  xleftmargin=14pt,
  breaklines=true,
  showstringspaces=false,
  tabsize=4,
  morekeywords={apply_chat_template,AutoTokenizer,AutoModelForCausalLM,model,tokenizer,messages,outputs,generate,decode},
  literate={<}{{\textless}}1 {>}{{\textgreater}}1 {|}{{\textbar}}1,
}

\lstdefinestyle{chatml}{
  backgroundcolor=\color{codeBg},
  basicstyle=\ttfamily\small,
  frame=l,
  framesep=4pt,
  rulecolor=\color{accentPurple},
  xleftmargin=14pt,
  breaklines=true,
  showstringspaces=false,
}

% ── Custom boxes ──────────────────────────────────────────────────────────────
\tcbset{
  notebox/.style={
    colback=noteYellow, colframe=noteYellowBorder,
    fonttitle=\bfseries\small, title=#1,
    left=6pt, right=6pt, top=4pt, bottom=4pt,
    breakable, arc=3pt
  },
  alertbox/.style={
    colback=lightRed, colframe=alertRed,
    fonttitle=\bfseries\small\color{white},
    coltitle=white, colbacktitle=alertRed,
    title=#1, left=6pt, right=6pt, top=4pt, bottom=4pt,
    breakable, arc=3pt
  },
  tipbox/.style={
    colback=lightTeal, colframe=teal,
    fonttitle=\bfseries\small\color{white},
    coltitle=white, colbacktitle=teal,
    title=#1, left=6pt, right=6pt, top=4pt, bottom=4pt,
    breakable, arc=3pt
  },
  defbox/.style={
    colback=lightPurple, colframe=accentPurple,
    left=6pt, right=6pt, top=4pt, bottom=4pt,
    breakable, arc=3pt
  },
}

% ── Hyperref ──────────────────────────────────────────────────────────────────
\hypersetup{
  colorlinks=true,
  linkcolor=primaryBlue,
  urlcolor=teal,
  pdftitle={Chat Templates in Transformers},
  pdfauthor={AI Engineering Notes},
}

% ── Helpers ───────────────────────────────────────────────────────────────────
\newcommand{\kw}[1]{\texttt{\color{primaryBlue}\bfseries #1}}
\newcommand{\val}[1]{\texttt{\color{codeGreen}#1}}
\newcommand{\tok}[1]{\texttt{\color{accentPurple}#1}}
\newcommand{\api}[1]{\texttt{\color{teal}#1}}

% ─────────────────────────────────────────────────────────────────────────────
\begin{document}

% ── Title block ───────────────────────────────────────────────────────────────
\begin{tcolorbox}[
  colback=headerBg, colframe=headerBg,
  coltext=white, arc=6pt,
  left=16pt, right=16pt, top=12pt, bottom=12pt
]
  {\LARGE\bfseries Chat Templates in Transformers}\\[6pt]
  {\large\color{lightBlue} \texttt{apply\_chat\_template()} --- Complete Engineering Notes}\\[4pt]
  {\small\color{borderGray} HuggingFace Transformers $\cdot$ ChatML $\cdot$ SmolLM2 $\cdot$ Qwen2.5 $\cdot$ Llama-3 $\cdot$ Production Patterns}
\end{tcolorbox}

\vspace{6pt}

% ── TOC ───────────────────────────────────────────────────────────────────────
{\small\color{mutedText}
\textbf{Contents:}
\S1 What \& Why \quad
\S2 API Reference \quad
\S3 ChatML Format \quad
\S4 Dry-Run Walkthrough \quad
\S5 Cross-Model Comparison \quad
\S6 Multimodal \quad
\S7 Production Patterns \quad
\S8 Pro Tips
}

\vspace{4pt}
\noindent\color{borderGray}\rule{\linewidth}{0.4pt}
\vspace{2pt}

% ═════════════════════════════════════════════════════════════════════════════
\section{What \& Why}
% ═════════════════════════════════════════════════════════════════════════════

Modern instruction-tuned LLMs are \textbf{not} trained on raw text. They are fine-tuned on structured conversations with role markers, special tokens, and precise formatting. Feeding unformatted text to an instruction model produces garbage output even if the content is correct.

\textbf{Chat templates} solve this by shipping a \textbf{Jinja2 template string} inside every model's \texttt{tokenizer\_config.json}. The \api{apply\_chat\_template()} method renders your \texttt{messages} list into the exact prompt format the model was trained on --- automatically, across any model family.

\begin{tcolorbox}[defbox]
  \textbf{Why not just write format strings manually?}\\
  Each model family uses a different format (ChatML, Llama-3 Header tokens, Mistral \texttt{[INST]}, Zephyr \texttt{<|system|>}). Manual code breaks across models and must be updated with every new release. Templates are self-contained and versioned with the tokenizer.
\end{tcolorbox}

% ═════════════════════════════════════════════════════════════════════════════
\section{API Reference --- \texttt{apply\_chat\_template()}}
% ═════════════════════════════════════════════════════════════════════════════

\subsection{Function Signature}

\begin{lstlisting}[style=pycode]
output = tokenizer.apply_chat_template(
    messages,                      # List[Dict[str, str]]  (required)
    tokenize=True,                 # False -> str  |  True -> BatchEncoding
    return_tensors="pt",           # "pt" | "tf" | "np" (only when tokenize=True)
    add_generation_prompt=True,    # appends <|im_start|>assistant\n
    return_dict=True,              # returns BatchEncoding with masks
    chat_template=None,            # override Jinja2 string (optional)
    add_special_tokens=True,       # add BOS/EOS (set False for non-chat)
)
\end{lstlisting}

\subsection{Parameter Deep-Dive}

\renewcommand{\arraystretch}{1.3}
\begin{tabularx}{\linewidth}{>{\ttfamily\color{primaryBlue}}p{5.2cm} X}
\toprule
\normalfont\textbf{Parameter} & \textbf{Behaviour \& Notes} \\
\midrule
messages & List of role dicts: \texttt{\{"role": "system|user|assistant", "content": "..."\}}. Must contain at least one user message. \\
\addlinespace
tokenize=False & Returns a formatted \textbf{string}. Use this for debugging --- always inspect the string before running inference. \\
\addlinespace
tokenize=True & Returns a \textbf{BatchEncoding} containing \texttt{.input\_ids} (and masks if \texttt{return\_dict=True}). Shape: \texttt{[1, seq\_len]}. \\
\addlinespace
return\_tensors="pt" & PyTorch tensors. Also accepts \texttt{"tf"} (TensorFlow) and \texttt{"np"} (NumPy). Only active when \texttt{tokenize=True}. \\
\addlinespace
add\_generation\_prompt=True & \textbf{Critical for inference.} Appends the assistant role-starter token(s) at the end so the model knows to generate a new turn. Omit only during fine-tuning data preparation. \\
\addlinespace
return\_dict=True & Returns full \texttt{BatchEncoding} with \texttt{input\_ids} \emph{and} \texttt{attention\_mask}. Required for batched inference with padding. \\
\addlinespace
chat\_template & Override the model's Jinja2 template inline without modifying the tokenizer object. \\
\bottomrule
\end{tabularx}

\subsection{Output Types}

\begin{tcolorbox}[notebox={Mode 1 --- Debug / Inspect}]
\begin{lstlisting}[style=pycode, numbers=none]
formatted = tokenizer.apply_chat_template(
    messages, tokenize=False, add_generation_prompt=True
)
print(formatted)   # human-readable string -> verify format before inference
\end{lstlisting}
\end{tcolorbox}

\begin{tcolorbox}[tipbox={Mode 2 --- Inference Ready}]
\begin{lstlisting}[style=pycode, numbers=none]
inputs = tokenizer.apply_chat_template(
    messages, tokenize=True, return_tensors="pt",
    add_generation_prompt=True, return_dict=True
)
# inputs.input_ids   -> torch.Tensor  [1, seq_len]
# inputs.attention_mask -> torch.Tensor  [1, seq_len]
outputs = model.generate(**inputs, max_new_tokens=200)
response = tokenizer.decode(
    outputs[0][inputs.input_ids.shape[-1]:],
    skip_special_tokens=True
)
\end{lstlisting}
\end{tcolorbox}

% ═════════════════════════════════════════════════════════════════════════════
\section{ChatML Format --- The Dominant Standard}
% ═════════════════════════════════════════════════════════════════════════════

Most open-weight models --- SmolLM2, Qwen2.5, Llama-3.x, Mistral-v3, Phi-3, Gemma-2 --- use the \textbf{ChatML} format:

\begin{lstlisting}[style=chatml]
<|im_start|>role
content<|im_end|>
\end{lstlisting}

\textbf{Special tokens in ChatML:}

\begin{itemize}[noitemsep, topsep=2pt]
  \item \tok{<|im\_start|>} --- Turn opener. Precedes every role block.
  \item \tok{<|im\_end|>} --- Turn closer. Signals end of this role's content (often doubles as EOS).
  \item BOS token --- Model-family specific (see \S5).
  \item Generation prompt --- \tok{<|im\_start|>assistant\textbackslash n} appended by \texttt{add\_generation\_prompt=True}.
\end{itemize}

\subsection{Other Formats (for reference)}

\renewcommand{\arraystretch}{1.3}
\begin{tabularx}{\linewidth}{>{\bfseries}p{3.2cm} X p{3.8cm}}
\toprule
Model Family & Format Pattern & Opener/Closer \\
\midrule
ChatML (majority) & \texttt{<|im\_start|>role\textbackslash ncontent<|im\_end|>} & \tok{<|im\_start|>} / \tok{<|im\_end|>} \\
Llama-3.x & Header-based with \texttt{<|start\_header\_id|>} & \texttt{<|eot\_id|>} \\
Mistral (old) & \texttt{[INST] user [/INST] assistant} & \texttt{[INST]} / \texttt{[/INST]} \\
Zephyr & \texttt{<|system|>\textbackslash ncontent</s>} & \texttt{<|system|>} / \texttt{</s>} \\
\bottomrule
\end{tabularx}

% ═════════════════════════════════════════════════════════════════════════════
\section{Dry-Run Walkthrough}
% ═════════════════════════════════════════════════════════════════════════════

Given the following \texttt{messages} list:

\begin{lstlisting}[style=pycode]
messages = [
    {"role": "system",    "content": "You are a helpful assistant."},
    {"role": "user",      "content": "What is the capital of France?"},
    {"role": "assistant", "content": "The capital of France is Paris."},
    {"role": "user",      "content": "And what about Germany?"},
]
\end{lstlisting}

\noindent\textbf{Step 1 --- Jinja2 renders each dict into a ChatML block:}

\begin{lstlisting}[style=chatml]
<|im_start|>system
You are a helpful assistant.<|im_end|>
<|im_start|>user
What is the capital of France?<|im_end|>
<|im_start|>assistant
The capital of France is Paris.<|im_end|>
<|im_start|>user
And what about Germany?<|im_end|>
\end{lstlisting}

\noindent\textbf{Step 2 --- \texttt{add\_generation\_prompt=True} appends the assistant starter:}

\begin{lstlisting}[style=chatml]
<|im_start|>assistant
^--- model starts generating here
\end{lstlisting}

\noindent\textbf{Step 3 --- Tokenization} (\texttt{tokenize=True, return\_tensors="pt"}):

\begin{lstlisting}[style=chatml]
SmolLM2  : shape [1, 55]  first_ids=[1, 9690, 273, ...]
Qwen2.5  : shape [1, N]   first_ids=[151644, 8948, 198, ...]
                                       ^--- <|im_start|> id differs!
\end{lstlisting}

% ═════════════════════════════════════════════════════════════════════════════
\section{Cross-Model Token Comparison}
% ═════════════════════════════════════════════════════════════════════════════

Same ChatML markers \emph{visually} --- completely different integer IDs because each model has its own vocabulary.

\renewcommand{\arraystretch}{1.4}
\begin{tabularx}{\linewidth}{>{\bfseries}l >{\ttfamily\color{primaryBlue}}X >{\ttfamily\color{teal}}X}
\toprule
\normalfont\textbf{Token} & SmolLM2-1.7B-Instruct & Qwen2.5-0.5B-Instruct \\
\midrule
BOS             & \textless s\textgreater{}             & \textless|endoftext|\textgreater{} \\
EOS             & \textless/s\textgreater{}            & \textless|im\_end|\textgreater{} \\
PAD             & \textless unk\textgreater{}          & \textless|endoftext|\textgreater{} \\
UNK             & \textless unk\textgreater{}          & None \\
\textless|im\_start|\textgreater{} id & 1                   & 151644 \\
Vocab size      & $\sim$49152                         & $\sim$151936 \\
\bottomrule
\end{tabularx}

\begin{tcolorbox}[alertbox={Critical Gotcha}]
Always load the \textbf{per-model tokenizer}. Never share a tokenizer between model families. The same \tok{<|im\_start|>} string maps to different integer IDs --- mixing tokenizers silently corrupts model input and produces hallucinated/broken output.
\end{tcolorbox}

\subsection{Llama-3.x Specific: Date Injection}

Llama-3.x templates automatically inject a \texttt{Today Date} field into the system prompt. This makes inputs \textbf{non-deterministic} by default --- a problem for caching:

\begin{lstlisting}[style=pycode]
# Override the date for deterministic / cached prompts
tokenizer.apply_chat_template(
    messages, tokenize=False,
    date_string="23 Jan 2025"   # pin the date
)
# Or strip it from the rendered string if you cannot avoid it:
prompt = prompt.replace("Cutting Knowledge Date: December 2023\nToday Date: 26 Jul 2024\n\n", "")
\end{lstlisting}

% ═════════════════════════════════════════════════════════════════════════════
\section{Multimodal Extension (2025)}
% ═════════════════════════════════════════════════════════════════════════════

Multimodal models (LLaVA, Qwen-VL, Idefics, etc.) extend the template with a \texttt{type} key in each content block. The method lives on the \textbf{Processor} class (not the tokenizer):

\begin{lstlisting}[style=pycode]
from transformers import AutoProcessor, LlavaOnevisionForConditionalGeneration

processor = AutoProcessor.from_pretrained("llava-hf/llava-onevision-qwen2-0.5b-ov-hf")

messages = [{
    "role": "user",
    "content": [
        {"type": "image", "url": "http://example.com/cat.jpg"},
        {"type": "text",  "text": "What do you see?"},
    ]
}]

inputs = processor.apply_chat_template(
    messages,
    add_generation_prompt=True,
    tokenize=True,
    return_dict=True,
    return_tensors="pt"
)
# inputs now contains both input_ids AND pixel_values
outputs = model.generate(**inputs, max_new_tokens=100)
\end{lstlisting}

\noindent Supported \texttt{type} values: \val{"text"}, \val{"image"}, \val{"video"}.\\
For video, pass \texttt{video\_load\_backend} parameter to control the video loading framework.

% ═════════════════════════════════════════════════════════════════════════════
\section{Production Patterns}
% ═════════════════════════════════════════════════════════════════════════════

\subsection{Full Inference Loop (Single Turn)}

\begin{lstlisting}[style=pycode]
import torch
from transformers import AutoTokenizer, AutoModelForCausalLM

tokenizer = AutoTokenizer.from_pretrained("HuggingFaceTB/SmolLM2-1.7B-Instruct")
model = AutoModelForCausalLM.from_pretrained("HuggingFaceTB/SmolLM2-1.7B-Instruct",
                                              torch_dtype=torch.bfloat16, device_map="auto")
messages = [
    {"role": "system",  "content": "You are a helpful assistant."},
    {"role": "user",    "content": "Explain RAG in one sentence."},
]

inputs = tokenizer.apply_chat_template(
    messages, tokenize=True, return_tensors="pt",
    add_generation_prompt=True, return_dict=True
)
inputs = {k: v.to(model.device) for k, v in inputs.items()}

outputs = model.generate(**inputs, max_new_tokens=200, do_sample=True, temperature=0.7)
response = tokenizer.decode(
    outputs[0][inputs["input_ids"].shape[-1]:],
    skip_special_tokens=True
)
print(response)
\end{lstlisting}

\subsection{Multi-Turn Conversation Loop}

\begin{lstlisting}[style=pycode]
messages = [{"role": "system", "content": "You are a helpful assistant."}]

while True:
    user_input = input("You: ")
    messages.append({"role": "user", "content": user_input})

    inputs = tokenizer.apply_chat_template(
        messages, tokenize=True, return_tensors="pt",
        add_generation_prompt=True, return_dict=True
    )
    outputs = model.generate(**inputs, max_new_tokens=300)
    response = tokenizer.decode(
        outputs[0][inputs["input_ids"].shape[-1]:], skip_special_tokens=True
    )
    # Append assistant turn BACK into messages to maintain context
    messages.append({"role": "assistant", "content": response})
    print(f"Model: {response}")
\end{lstlisting}

\subsection{Custom / Override Template}

\begin{lstlisting}[style=pycode]
# Option A: mutate the tokenizer object (persists for all calls)
tokenizer.chat_template = your_jinja2_string

# Option B: pass inline per-call (no mutation)
output = tokenizer.apply_chat_template(
    messages, tokenize=False, chat_template=your_jinja2_string
)

# Option C: edit tokenizer_config.json on disk before loading
# Useful when the model's template is buggy (e.g. wrong date injection)
\end{lstlisting}

\subsection{Fine-Tuning Data Preparation}

\begin{lstlisting}[style=pycode]
# During fine-tuning: include the assistant response in messages
# and set add_generation_prompt=False
messages = [
    {"role": "user",      "content": "What is the capital of France?"},
    {"role": "assistant", "content": "The capital of France is Paris."},
]

# Transformers >= 4.42: get label mask for SFT (ignore prompt tokens)
result = tokenizer.apply_chat_template(
    messages,
    tokenize=True,
    return_tensors="pt",
    add_generation_prompt=False,
    return_dict=True,
    return_assistant_tokens_mask=True,  # mask non-assistant tokens from loss
)
# result["assistant_masks"] -> 1 where loss should be computed
\end{lstlisting}

\subsection{LangChain Integration}

When using \texttt{HuggingFacePipeline} in LangChain, \api{apply\_chat\_template()} is called \textbf{automatically} under the hood when the model is a chat/instruction model. You do not need to call it manually --- just pass your messages and LangChain handles the rest.

% ═════════════════════════════════════════════════════════════════════════════
\section{Pro Tips for AI Engineering}
% ═════════════════════════════════════════════════════════════════════════════

\begin{tcolorbox}[notebox={Debug-First Workflow}]
\textbf{Always} inspect the formatted string before feeding tensors to the model:
\begin{lstlisting}[style=pycode, numbers=none]
print(tokenizer.apply_chat_template(messages, tokenize=False, add_generation_prompt=True))
\end{lstlisting}
This costs nothing and catches format bugs instantly.
\end{tcolorbox}

\begin{itemize}[noitemsep, topsep=4pt, leftmargin=1.2em]

  \item \textbf{\texttt{add\_generation\_prompt=True} is mandatory for inference.}
    Omitting it causes the model to try to continue the conversation rather than generate a fresh assistant turn. Symptoms: repeated content, wrong role, or immediate EOS.

  \item \textbf{Use \texttt{return\_dict=True} for batched inference.}
    Without the \texttt{attention\_mask} returned by \texttt{return\_dict=True}, padding in batches is ignored and generation quality degrades.

  \item \textbf{\texttt{return\_assistant\_tokens\_mask=True} for SFT.}
    Available from Transformers $\ge$ 4.42. Automatically generates the loss mask that suppresses gradients on prompt tokens --- replaces manual label masking logic.

  \item \textbf{Never mix tokenizers across model families.}
    The same special token string maps to different integer IDs per vocab. Load \texttt{AutoTokenizer.from\_pretrained(model\_name)} for each model individually.

  \item \textbf{Messages list must contain at least one user message.}
    A system-only \texttt{messages} list raises a Jinja2 rendering error in most templates.

  \item \textbf{Slice outputs to get only the new tokens:}
    \texttt{outputs[0][inputs["input\_ids"].shape[-1]:]} strips the echoed prompt from the generated sequence.

  \item \textbf{For RAG pipelines:} inject retrieved context into the \texttt{system} role content or prepend it to the first \texttt{user} message. The template handles token budget automatically across the full conversation window.

  \item \textbf{Customise templates carefully.}
    If you edit \texttt{tokenizer.chat\_template}, test with \texttt{tokenize=False} first, verify the rendered string against the model's expected format, then re-enable tokenisation.

\end{itemize}

% ── Quick reference card ──────────────────────────────────────────────────────
\vspace{6pt}
\noindent\color{borderGray}\rule{\linewidth}{0.4pt}
\vspace{2pt}

\begin{tcolorbox}[
  colback=codeBg, colframe=borderGray,
  title={\color{darkText}\bfseries Quick Reference --- apply\_chat\_template()},
  fonttitle=\bfseries
]
\renewcommand{\arraystretch}{1.25}
\begin{tabularx}{\linewidth}{>{\ttfamily\footnotesize\color{primaryBlue}}p{5.5cm} >{\footnotesize}X}
\toprule
\normalfont\textbf{Call} & \textbf{Returns / Use case} \\
\midrule
apply\_chat\_template(m, tokenize=False) & Formatted string --- debug \\
apply\_chat\_template(m, tokenize=True, return\_tensors="pt") & Bare tensor \texttt{input\_ids} \\
apply\_chat\_template(m, return\_dict=True) & BatchEncoding with masks --- batching \\
apply\_chat\_template(m, add\_generation\_prompt=True) & Add assistant starter --- inference \\
apply\_chat\_template(m, add\_generation\_prompt=False) & No suffix --- fine-tuning data \\
apply\_chat\_template(m, return\_assistant\_tokens\_mask=True) & SFT loss mask (HF $\ge$ 4.42) \\
apply\_chat\_template(m, date\_string="...") & Override Llama-3 date injection \\
\bottomrule
\end{tabularx}
\end{tcolorbox}

\vspace{4pt}
\begin{center}
  \small\color{mutedText}
  Source: HuggingFace Transformers Docs $\cdot$
  \href{https://huggingface.co/docs/transformers/en/chat_templating}{docs/transformers/chat\_templating} $\cdot$
  Compiled April 2026
\end{center}

\end{document}